\documentclass{article}
\usepackage{amsmath,amssymb,mathtools}
\usepackage[margin=1in]{geometry}
\allowdisplaybreaks
\usepackage[UTF8]{ctex}
\begin{document}

\section*{算法1：注重隐私的动态个性化定价}

\subsection*{收入目标与定价规则}
零售商旨在实现最大化
\[
\mathbb E\!\left[\sum_{t=1}^{T}p_ty_t\right],
\qquad p_t\in[0,1],\quad t=1,\ldots,T.
\]
所选的策略是算法1，但这并不意味着其价格能够精确实现这一累积目标的最大化。其价格选择的行是：
\[
p_n=\arg\max_{p\in[0,1]}\min\left\{1,
 p f\!\left(\phi_n^{\top}\widehat\theta^{p}\right)
 +\gamma\sqrt{\phi_n^{\top}(\Lambda^{p})^{-1}\phi_n}\right\},
\qquad \phi_n=\phi(x_n,p_n),
\quad n=T_0+1,\ldots,T.
\]
这一行中列出的状态与释放情况如下所述。该设置为一个单一的产品，拥有无限的物品库存（pp.~2--3）。

\paragraph{印刷版价格标注。}
算法1的第10行中，$\phi_n$被定义为使用$p_n$，而不是优化变量$p$。这种打印出的依赖关系被保留下来，不会被$\phi(x_n,p)$所替代。

\subsection*{需求、信息以及固定模型假设}
我们保留了第3节中的抽象指数族模型（pp.~9--10），而不对其中的示例性可能性进行专门化处理。在$t=1,\ldots,T$、$x_t$这些时期，分别代表到达的客户的个人信息、$p_t$代表提供的价格，而$y_t\in\mathbb R$则代表实际的需求。已知的特征映射和未知的固定模型则保持不变。
\[
\phi_t=\phi(x_t,p_t)\in\mathbb R^d,
\qquad \theta^*\in\mathbb R^d.
\]
在这篇摘要中，并未明确指定$x_t$的特定维度或坐标域。分布家族组件$\zeta,m(\cdot),h(\cdot)$仍然保持抽象状态。印刷需求模型是
\begin{equation}
\Pr[y_t=y\mid p_t,x_t,\theta^*]
=\exp\!\left\{\zeta\left(y\phi_t^{\top}\theta^*
-m(\phi_t^{\top}\theta^*)\right)+h(y)\right\}.
\tag{1}
\end{equation}
其条件均值和噪声表示分别为：
\[
f(\phi_t^{\top}\theta^*):=m'(\phi_t^{\top}\theta^*)
=\mathbb E[y_t\mid p_t,x_t,\theta^*],
\]
\begin{equation}
y_t=f(\phi_t^{\top}\theta^*)+\xi_t.
\tag{2}
\end{equation}
第 (2) 段指出，$\phi_t=\phi(x_t,p_t)$和$\xi_t$是独立的随机变量，它们满足……
\[
\mathbb E[\xi_t\mid p_t,x_t]=0.
\]
历史以及后续的条件性居中操作如下：
\[
\mathcal F_t=\{(x_\tau,y_\tau,p_\tau)\}_{\tau=1}^{t},
\qquad
\mathbb E[\xi_t\mid x_t,p_t,\mathcal F_{t-1}]=0.
\]
在本文中使用的通用设置中，$x_1,\ldots,x_T$是在定价开始之前任意选择的（即在“ oblivious-adversary”设置下）。该策略基于$\mathcal F_{t-1}$和$x_t$来选择$p_t$，之后，需求将按照公式(2)来得到满足。

在整篇论文中，来源方提出了以下模型假设。这些假设涉及固定输入和需求过程，而不是对估计器的额外优化约束。
\begin{align*}
& B_Y<\infty,\qquad |y_t|\le B_Y
&&\text{for every period }t\text{ in every database }D;\\
&\|\phi(x,p)\|_2\le1,\qquad \|\theta^*\|_2\le1
&&\text{for all }x,p;\\
&\mathbb E[\xi_t\mid x_t,p_t,\mathcal F_{t-1}]=0,\\
&\text{there exists }s<\infty\text{ such that }
\mathbb E[e^{\lambda\xi_t}\mid x_t,p_t,\mathcal F_{t-1}]
\le e^{\lambda^2s^2/2}
&&\text{for all }\lambda\in\mathbb R;\\
&f=m':\mathbb R\longrightarrow[0,1],
\qquad f\text{ continuously differentiable and strictly increasing};\\
&1\le K<\infty,\qquad K^{-1}\le f'(z)\le K
&&\text{for all }|z|\le2;\\
&1\le G<\infty,\qquad G^{-1}\le\zeta\le G.
\end{align*}
这里的$B_Y,s,K,G$是模型参数。在算法3所使用的似然表示法中，公式(1)中的函数，在候选参数$\theta$处的值，就是……
\[
p(y_t\mid\phi_t,\theta)
=\exp\!\left\{\zeta\left(y_t\phi_t^{\top}\theta
-m(\phi_t^{\top}\theta)\right)+h(y_t)\right\}.
\]
似然符号$p(\,\cdot\mid\cdot\,)$与标量价格$p$是不同的。这里使用的是第3节的内容，而不是p.~25上所写的另一种表述方式。

\subsection*{算法1的输入、初始化以及顺序更新}
在 p.~20 上的固定算法输入是：
\[
\varepsilon_1,\delta_1,\varepsilon_2,\delta_2>0,
\qquad \rho\ge1,\qquad \gamma>0,
\]
结合纯探索周期的数量$T_0$和最大调用次数\textsc{PrivateMLE}调用$D_\infty$。输出结果为$p_1,p_2,\ldots,p_T$。量$\widehat\theta^p$、$\Lambda^p$和$D_{\mathrm{MLE}}$是算法维护的状态，而非固定输入；上标$p$表示私有化量。

初始化，与第3行中的操作完全相同。
\[
\delta'_2=\frac{\delta_2}{2D_\infty},
\qquad
\varepsilon'_2\leftarrow
\frac{\varepsilon_2}{2\sqrt{2D_\infty\ln(1/\delta'_2)}},
\qquad
\Lambda^p=\rho I_d,
\qquad \widehat\theta^p=0,
\qquad D_{\mathrm{MLE}}=0.
\]
对于前$T_0$个周期，以均匀随机的方式提供价格：
\[
p_t\sim U[0,1],\qquad t=1,\ldots,T_0.
\]
对于每一个$n=T_0+1,\ldots,T$，首先获取
\[
\Sigma^p_{\mathrm{used},n}\leftarrow R_{n-1},\quad R_j:=\text{the output of iteration }j\text{ of the single horizon-}T\text{ covariance tree},
\qquad \Lambda_n^p=\Sigma^p_{\mathrm{used},n}+\rho I_d.
\]
如果
\[
\det(\Lambda_n^p)>2\det(\Lambda^p)
\quad\text{and}\quad D_{\mathrm{MLE}}<D_\infty,
\]
执行这些任务/操作。
\[
\widehat\theta^p\leftarrow
\textsc{PrivateMLE}(n,\rho,\varepsilon'_2,\delta'_2),
\qquad \Lambda^p\leftarrow\Lambda_n^p,
\qquad D_{\mathrm{MLE}}\leftarrow D_{\mathrm{MLE}}+1.
\]
否则，请保留这些状态不变。然后，按照上述价格选择行所给出的价格来报价。在这里，$I_d$是$d$维的单位矩阵，而与协方差相关的矩阵的维数为$d\times d$。

注意 4（pp.~20–21）允许省略强制探索步骤，使用$T_0=0$来表示。对于此处复现的参数化算法1而言，没有针对特定定理的调整或随机上下文假设。

\subsection*{算法 2: \textsc{PrivateCov}}
算法 2（p.~21）的签名为：
\[
\textsc{PrivateCov}(T,\varepsilon,\delta)
\quad\text{returns}\quad\Sigma_1^p,\ldots,\Sigma_{T-1}^p.
\]
其局部参数和初始状态为：
\begin{align*}
\delta'&\leftarrow\frac{\delta}{2\lceil\log_2T\rceil},
&\varepsilon'&\leftarrow
\frac{\varepsilon}{2\lceil\log_2T\rceil\ln(1/\delta')},\\
\sigma_{\varepsilon',\delta'}^2
&=\frac{2\ln(1.25/\delta')}{(\varepsilon')^2},
&m&=\lceil\log_2T\rceil,\\
\Sigma(\ell)&=\widehat\Sigma(\ell)=0,
&&\ell=0,\ldots,m-1.
\end{align*}
这里的局部整数$m$表示二叉树的深度，它与分布族函数$m(\cdot)$是不同的。

对于$n=1,2,\ldots,T-1$，将$n$表示为二进制形式，并确定其最低有效非零位：
\[
n=\sum_{\ell=0}^{m-1}b_n(\ell)2^\ell,
\qquad b_n(\ell)\in\{0,1\},
\qquad \ell_n\leftarrow\min\{\ell:b_n(\ell)=1\}.
\]
以打印时的顺序来更新和重置：
\[
\Sigma(\ell_n)\leftarrow
\phi_n\phi_n^{\top}+\sum_{\ell<\ell_n}\Sigma(\ell),
\qquad
\Sigma(\ell)\leftarrow\widehat\Sigma(\ell)\leftarrow0
\quad\text{for all }\ell<\ell_n.
\]
校准噪声并释放：
\[
\widehat\Sigma(\ell_n)\leftarrow\Sigma(\ell_n)+W^n,
\qquad
W^n_{ij}=W^n_{ji}\overset{\mathrm{i.i.d.}}{\sim}
\mathcal N(0,\sigma_{\varepsilon',\delta'}^2),
\]
\[
\Sigma_n^p=\sum_{\ell=0}^{m-1}b_n(\ell)\widehat\Sigma(\ell).
\]
来源中将$W^n$描述为一个对称的高斯矩阵；其显示的对称性和独立同分布的特性被保留下来，没有添加任何采样规则。第6.1节（p.~22）定义了相应的未受扰动样本协方差。
\[
\Sigma_n=\sum_{t=1}^{n}\phi_t\phi_t^{\top},
\qquad n=1,\ldots,T-1.
\]

\paragraph{印刷发布时间冲突。}
第6.1节中，位于~21处的内容指出：$\textsc{PrivateCov}(n,\varepsilon,\delta)$是迭代$n-1$结束时释放的$\Sigma_n^p$矩阵。而算法2则将其在迭代$n$时的释放过程标记为$\Sigma_n^p$，并在该迭代中引入$\phi_n\phi_n^{\top}$。算法1在选择$p_n$之前会调用$\textsc{PrivateCov}(n,\varepsilon_1,\delta_1)$，而第5.1节则要求，在时间$n$使用的释放函数只能访问更早的观测结果。这里保留了打印出来的冲突信息，作为来源记录。所选版本使用了第6.1节中的句子：用于决策的协方差$n$，就是迭代$n-1$结束时释放的那个协方差。该树结构保留了打印出来的迭代索引，并且每个历史行都被使用了一遍。特别地，算法2并没有为允许的$T_0=0$情况指定一个零迭代时的释放过程。

\subsection*{算法 3: \textsc{PrivateMLE}}
在请求的通话中，
\[
\textsc{PrivateMLE}(n,\rho,\varepsilon,\delta)
\quad\text{returns}\quad\widehat\theta_n^p.
\]
算法1提供$(\varepsilon,\delta)=(\varepsilon'_2,\delta'_2)$。算法3（p.~24）首先进行赋值操作。
\[
B_1\leftarrow(B_Y+1)G,
\qquad B_2\leftarrow KG,
\qquad \rho\leftarrow\max\{\rho,2B_2/\varepsilon\},
\]
\[
\nu_{\varepsilon,\delta}^{2}\leftarrow
\frac{B_1^2\bigl(8\ln(2/\delta)+4\varepsilon\bigr)}{\varepsilon^2},
\qquad
w\sim\mathcal N(0,\nu_{\varepsilon,\delta}^{2}I_d).
\]
然后返回
\[
\widehat\theta_n^p
=\arg\min_{\|\theta\|_2\le2}
\left\{\left(\sum_{t<n}-\ln p(y_t\mid\phi_t,\theta)\right)
+\frac{\rho}{2}\|\theta\|_2^2+w^{\top}\theta\right\}.
\]
这里，$\theta\in\mathbb R^d$是估计变量，$w\in\mathbb R^d$是采样扰动，而$t<n$则涵盖了呼叫之前的观测周期。半径二估计约束与固定模型假设$\|\theta^*\|_2\le1$是不同的。$\rho$的重新分配在算法3中得到了体现；算法1中的协方差状态没有进行额外的重新分配。

\paragraph{打印版隐私分配差异。}
算法1中的分配方式将因子$2$保留在$\varepsilon'_2$的分母中，且位于平方根之外。而附录3中关于p.~23的分配方式则省略了该因子。其替代分配方式并未被纳入算法1中。

\subsection*{政策附带的隐私定义}
对于隐私定义（pp.~12--14），请写入
\[
D=\{o_t\}_{t=1}^{T},\qquad o_t=(x_t,y_t),
\qquad D'=\{o'_t\}_{t=1}^{T}.
\]
在时间$t$时，存在差异的邻近数据库能够满足以下条件：
\[
o_t\ne o'_t,\qquad o_\tau=o'_\tau\quad\text{for all }\tau\ne t.
\]
对于$\varepsilon,\delta>0$，要实现与$(\varepsilon,\delta)$不同的差分隐私，就需要针对每一对这样的元素，以及每一个可测量的集合$\mathcal P_{>t}$，采取相应的措施。
\begin{equation}
\Pr[p_{t+1},\ldots,p_T\in\mathcal P_{>t}\mid\pi,D]
\le e^{\varepsilon}
\Pr[p_{t+1},\ldots,p_T\in\mathcal P_{>t}\mid\pi,D']+\delta.
\tag{3}
\end{equation}
这些私有释放器仅使用历史数据，而价格优化器则使用它们的输出结果以及当前的$x_t$，而不是直接使用历史敏感数据（参见~17–18页）。所声称的复合隐私级别为：
\[
\varepsilon=\varepsilon_1+\varepsilon_2,
\qquad \delta=\delta_1+\delta_2.
\]
此隐私定义是无需模型的；上述需求假设用于效用分析，而非用于定义隐私。第~23页的推论4将所得价格序列称为“差分隐私的”，但没有使用“预期”一词，尽管其引用的命题2实际上是在讨论预期对差分隐私的影响。此处保留的定义是（3），而不是更强的全价格序列要求。

\paragraph{来源范围。}
操作性的表述是参数化的算法1及其与算法2-3的依赖关系，同时还包括第3节中的需求族。高斯分布和逻辑斯蒂分布的相关示例、特定定理的参数选择、后悔分析不等式以及仿真设置等，均不属于额外的操作性约束。本文将证明内容附在在线补充材料中；已发布的电子配套文件EC.2则使用了从时期$t-1$开始的历史总结数据，其价格为$t$。EC.3.1则采用了基于单一视界的树深度方法。为pp.~3、20、21和24提供的页面图片，支持了上述目标和算法描述的内容。


\paragraph{对存档的基准输入进行选择。}
对于所提供的十个标量输入，\(T_0\geq2\)在每次适应性决策之前都会出现一次，包括最后一个决策。实现方式选择打印价格选择器中的最大化值。每个案例都有与价格无关的\(\phi(x,p)=(1+x)/2\)和逻辑均值。在初始化时，标量接受矩阵为正数；严格的加倍门限确保了正性。因此，在每个实现的历史中，上限乐观得分随着提供价格的上升而递增，而价格为一时则是一个最大化选择。统一的探索性报价仍然是必需的。这些实例事实并未指定与价格相关的或多维情况。只有未来的公开报价才会被纳入预期隐私声明中。存档的播种诊断路径、噪声发布、购买和估计值都支持复制，但这些内容并非带有隐私保证的公开发布。
\end{document}
